% §5 The trajectory atlas
The trajectory atlas renders the coverage computation as one behavior plane: real walks as
blue dots, simulated walks as orange diamonds, uncovered cells shaded red with
nearest-exemplar tooltips, outcome classes readable by position (and by per-class views as
small multiples).
It is simultaneously the paper's worked example and the artifact's operational interface.
Figure~\ref{fig:atlas} shows the schema.

\begin{figure}[t]
\centering
\includegraphics[width=0.9\textwidth]{fig_atlas}
\caption{The trajectory atlas, rendered from the artifact's coverage computation (same
deterministic run as Figure~\ref{fig:coverage}: 2 real / 1 simulated fixture walks,
rank-deficient basis at $n{=}2$ --- a smoke test of runnability, §\ref{sec:evaluation}).
Red cells are the 153 occupied-but-uncovered regions (top-3 by local real density outlined);
each carries its nearest real exemplar.
A concentrated sim cloud with a small radius (over-cooperative collapse) and a scattered sim
cloud with uncovered lobes is the typical deficient configuration.}
\label{fig:atlas}
\end{figure}

\subsection{Reading the atlas}
Three configurations recur.
\emph{Collapse}: the sim cloud concentrates on one real lobe (the over-cooperative user who
never abandons, never escalates --- cf.\ the Formalism Ceiling \citep{realUserSim2026} and the
human-likeness failures of \citep{mirrorbench2026}); coverage is high on the majority class
and near zero elsewhere.
\emph{Scatter}: the sim cloud spreads where real conversations never go --- off-distribution
generation that inflates the plane without measuring anything; the shared threshold of
Equation~\eqref{eq:threshold} keeps these points from manufacturing coverage, and
out-of-grid points are dropped rather than clamped.
\emph{Coverage with holes}: the sim cloud overlaps real mass but misses specific occupied
cells --- the most actionable configuration, because the holes come with named exemplars.
Cross-model comparison is legitimate by design: real logs and simulated runs may come from
different vendors and model generations, and the representation comparison rests on
cross-model convergence of layer representations \citep{layersSimilarDepths2025} and the
trajectory-geometry vocabulary \citep{trajectoryGeometry2026}.

\subsection{Worked example on shipped fixtures}
The open-source artifact ships fixtures in three dialects and regenerates every number below
from a single command.
The worked example uses 3 conversations across 3 dialects --- two real (one OTel GenAI export,
one OpenInference export) and one simulated (a LangSmith export) --- under the grammar feature
encoding, a $24 \times 24$ grid, $\varepsilon = 0.05$, bandwidth $\beta = 0.12$.
The real set occupies 256 grid cells; the 1-conversation simulated set covers 103 of them.

\begin{quote}
\small
\texttt{coverage: 40.2\% overall | per class: contained 45.3\%, handoff 35.2\%}
\end{quote}

\begin{table}[t]
\centering
\small
\begin{tabular}{@{}clllc@{}}
\toprule
\textbf{Rank} & \textbf{Grid cell} & \textbf{Local density} & \textbf{Nearest real exemplar (first user utterance)} & \textbf{Class} \\
\midrule
1 & (2,\,11)  & 0.86 & ``No puedo entrar en mi banca online'' & handoff \\
2 & (21,\,12) & 0.86 & ``Hola, quiero consultar el saldo de mi cuenta'' & contained \\
3 & (2,\,14)  & 0.54 & ``No puedo entrar en mi banca online'' & handoff \\
\bottomrule
\end{tabular}
\caption{Top-3 uncovered regions from the worked example (as reported by the shipped CLI).
Each names its nearest real exemplar --- the tests the simulated set is not running.
Exemplars are shown in the fixture's original language (Spanish); the pipeline is
language-agnostic by construction, with clarification/confirmation pattern tables as the only
language-coupled component.}
\label{tab:topk}
\end{table}

The interpretation is the paper's pitch in one line: a one-conversation simulated set covers
under half of the conversational-behavior regions that the real set occupies --- and the
uncovered regions are not an abstraction, they are named conversations with outcome classes
and cycle depths (Table~\ref{tab:topk}).
The numbers are fixture-scale by design; what they demonstrate is that the metric runs
end-to-end on standard exports and that its output format is directly actionable.
All numbers regenerate from the artifact's check suites (46 checks across extraction and
coverage, all passing at the time of writing) and the coverage CLI; the worked example is
reproducible from the repository README.

\subsection{Relation to the simulator program}
The atlas is where calibrated simulation is \emph{validated}.
The calibrated-simulation agenda --- response-level gap closure
\citep{distributionalGap2026}, calibrated user embeddings \citep{cue2026}, grounded user
simulation \citep{realUserSim2026} --- produces better simulators; region coverage is the
instrument that says whether the improvement reached the conversation structures that matter.
The secondary instruments for that validation are Wasserstein divergence of grammar-walk
statistics in the style of the Conversational Distribution Score \citep{cds2026} and the rank
correlation of Section~\ref{sec:evaluation}; region coverage is the headline, because it is
the only one of the three that returns actionable per-region output.
